% Copyright (c) 2026 David Immanuel Eidenas.
% SPDX-License-Identifier: CC-BY-4.0
\documentclass[11pt,a4paper]{amsart}
\usepackage[T1]{fontenc}
\usepackage{lmodern}
\usepackage{amsmath,amssymb,amsthm}
\usepackage{stmaryrd}
\usepackage[margin=30mm]{geometry}
\usepackage[hidelinks]{hyperref}

% Preserve the capitalisation entered in the author field.
\makeatletter
\let\publication@setauthors\@setauthors
\renewcommand{\@setauthors}{%
  \begingroup
  \let\MakeUppercase\@firstofone
  \publication@setauthors
  \endgroup
}
\makeatother

\hypersetup{
  pdftitle={Strict inequalities for products of filters},
  pdfauthor={David Immanuel Eidenas},
  pdfsubject={Preprint (not peer reviewed): an affirmative solution to Open Problem Garden problem 37339},
  pdfkeywords={filters, funcoids, reloids, oblique products, Lean}
}

\newtheorem{theorem}{Theorem}[section]
\newtheorem{lemma}[theorem]{Lemma}
\newtheorem{proposition}[theorem]{Proposition}
\newtheorem{corollary}[theorem]{Corollary}
\theoremstyle{definition}
\newtheorem{definition}[theorem]{Definition}
\theoremstyle{remark}
\newtheorem{remark}[theorem]{Remark}

\newcommand{\RLD}{\mathrm{RLD}}
\newcommand{\FCD}{\mathrm{FCD}}
\newcommand{\out}{\operatorname{out}}
\newcommand{\fil}{\operatorname{fil}}
\newcommand{\Up}{\operatorname{up}}
\newcommand{\Pos}{\operatorname{Pos}}
\newcommand{\N}{\mathbb{N}}
\newcommand{\R}{\mathbb{R}}
\newcommand{\cF}{\mathcal{F}}
\newcommand{\cG}{\mathcal{G}}
\newcommand{\cX}{\mathcal{X}}
\newcommand{\cY}{\mathcal{Y}}

\title{Strict inequalities for products of filters}
\author{David Immanuel Eidenas}
\thanks{Preprint---not peer reviewed. Licensed under
\href{https://creativecommons.org/licenses/by/4.0/}{CC BY 4.0}.}
\date{October 9, 2026}
\keywords{Filters, funcoids, reloids, oblique products, formal verification}

\begin{document}

\begin{abstract}
We answer affirmatively the two conjectures in Open Problem Garden problem
37339 concerning the second, oblique, and classical reloidal products of
filters. The cofinite filter $\Phi$ on $\N$ satisfies
$\Phi\times_F^{\RLD}\Phi=\bot\sqsubset\Phi\ltimes\Phi
\sqsubset\Phi\times^{\RLD}\Phi$.
The same strict chain holds for the proposed one-sided neighbourhood filter
$\Delta_+$ on $\R$. In both examples, two disjoint relations lying above the
funcoidal product force the second product to be the least reloid, whereas a
column characterisation makes the oblique product proper. The complement of
the diagonal separates the oblique and classical products. The accompanying
Lean development formalises both examples, including the real-number case.
\end{abstract}

\maketitle
\markboth{\shortauthors}{\shorttitle}

\section{The problem and conventions}

Porton's problem \cite{garden} asks whether there are filter objects
$\cF,\cG$ such that
\begin{equation}\label{eq:problem}
  \cF\times_F^{\RLD}\cG
  \sqsubset \cF\ltimes\cG
  \sqsubset \cF\times^{\RLD}\cG,
\end{equation}
and asks specifically about
$\cF=\cG=\Delta\sqcap\uparrow^{\R}(0,+\infty)$.
It also asks the weaker question of strictness of the first inequality alone.
We use the definitions in \cite{porton}, with the existential quantifier in
the oblique product as corrected by its author in 2016 \cite{erratum}.

A filter on a set $U$ is a nonempty, upward-closed family of subsets of $U$
closed under finite intersections. Improper filters are allowed; a filter is
\emph{proper} when it does not contain $\varnothing$.
Following the source, the order is
\[
  \cF\sqsubseteq\cG \quad\Longleftrightarrow\quad \cF\supseteq\cG,
\]
and $\sqsubset$ denotes its strict part. The least filter $\bot$ is
$\mathcal P(U)$; the principal filter $\uparrow A$ is
$\{S\subseteq U:A\subseteq S\}$.

For filters $\cF,\cG$ on the same set, write $\cF\mathrel{\#}\cG$ when
their meet is proper. Equivalently,
\begin{equation}\label{eq:mesh}
  \cF\mathrel{\#}\cG
  \quad\Longleftrightarrow\quad
  A\cap B\ne\varnothing \quad(A\in\cF,\ B\in\cG).
\end{equation}
Indeed, their meet is the filter generated by $\cF\cup\cG$, and every finite
intersection of members of that union reduces to an intersection of one
member of each filter. A set $X\subseteq U$ is \emph{$\cF$-positive} if
$\uparrow X\mathrel{\#}\cF$, that is, if it meets every member of $\cF$.
We denote the family of such sets by $\Pos(\cF)$.

A reloid from $U$ to $V$ is represented by a filter on $U\times V$. We write
$\Up(P)$ for its family of relations. Reloid order is therefore
$P\sqsubseteq Q$ if and only if $\Up(P)\supseteq\Up(Q)$.
For a family $J\subseteq\mathcal P(U\times V)$, put
\begin{equation}\label{eq:generated}
 \fil(J)=\left\{S\subseteq U\times V:
 \exists n\ge0,\ R_1,\ldots,R_n\in J,\
 R_1\cap\cdots\cap R_n\subseteq S\right\},
\end{equation}
with empty intersection $U\times V$. This is the family of the reloid meet
of the principal reloids generated by members of $J$.

For a relation $R\subseteq U\times V$, use the notation
\[
 R[X]=\{y\in V:\exists x\in X,\ (x,y)\in R\},
 \qquad C_y(R)=\{x\in U:(x,y)\in R\}.
\]

\section{The generating families}

The funcoidal product has the graph
\[
 \cX[\cF\times^{\FCD}\cG]\cY
 \quad\Longleftrightarrow\quad
 \cX\mathrel{\#}\cF\ \text{ and }\ \cY\mathrel{\#}\cG.
\]
For a relation $R$, the image of $\cX$ under its principal funcoid is the
filter generated by $\{R[A]:A\in\cX\}$, and
\[
 \cX[\uparrow^{\FCD}R]\cY
 \quad\Longleftrightarrow\quad
 \cY\mathrel{\#}\langle\uparrow^{\FCD}R\rangle\cX.
\]
Funcoids are ordered by inclusion of these graphs. Define
\begin{align*}
 I_2(\cF,\cG)&=\{R:\uparrow^{\FCD}R
                         \sqsupseteq\cF\times^{\FCD}\cG\},\\
 I_{\ltimes}(\cF,\cG)&=\{R:\exists B_0\in\cG,\
            \uparrow^{\FCD}R\sqsupseteq\cF\times^{\FCD}\uparrow B_0\},\\
 I_{\square}(\cF,\cG)&=\{A\times B:A\in\cF,\ B\in\cG\}.
\end{align*}
The definitions of the three products give
\begin{align}\label{eq:products}
 \Up(\cF\times_F^{\RLD}\cG)&=\fil(I_2(\cF,\cG)),\nonumber\\
 \Up(\cF\ltimes\cG)&=\fil(I_{\ltimes}(\cF,\cG)),\\
 \Up(\cF\times^{\RLD}\cG)&=\fil(I_{\square}(\cF,\cG)).\nonumber
\end{align}
In the first line, the defining operation
$(\RLD)_{\out}f=\bigsqcap^{\RLD}\Up(f)$ ranges over the
\emph{principal relations above} $f$.

\begin{lemma}\label{lem:dual}
For a filter $\cF$ on $U$ and a set $C\subseteq U$,
\[
 C\in\cF
 \quad\Longleftrightarrow\quad
 C\cap X\ne\varnothing\text{ for every }X\in\Pos(\cF).
\]
\end{lemma}
\begin{proof}
The forward implication is the definition of positivity. If $C\notin\cF$,
no member of $\cF$ is contained in $C$. Thus $U\setminus C$ is
$\cF$-positive and disjoint from $C$, proving the reverse implication.
\end{proof}

\begin{lemma}\label{lem:indices}
For filters $\cF$ on $U$, $\cG$ on $V$, and a relation $R\subseteq U\times V$:
\begin{align}
 R\in I_2(\cF,\cG)
 &\Longleftrightarrow
 R[X]\cap Y\ne\varnothing
 \quad(X\in\Pos(\cF),\ Y\in\Pos(\cG));\label{eq:i2}\\
 R\in I_{\ltimes}(\cF,\cG)
 &\Longleftrightarrow
 \exists B_0\in\cG\ \forall y\in B_0:\ C_y(R)\in\cF.\label{eq:ilt}
\end{align}
\end{lemma}
\begin{proof}
For \eqref{eq:i2}, necessity follows by testing the graph order on principal
filters $\uparrow X,\uparrow Y$. Conversely, if
$\cX\mathrel{\#}\cF$ and $\cY\mathrel{\#}\cG$, every $A\in\cX$ is
$\cF$-positive and every $B\in\cY$ is $\cG$-positive. Hence $R[A]\cap B$
is nonempty for all such members. The images $R[A]$ form a filter base:
$R[A_1\cap A_2]\subseteq R[A_1]\cap R[A_2]$. Therefore $\cY$ meshes with
the image filter, giving the required graph inclusion.

For a fixed $B_0$, the same graph condition with second factor
$\uparrow B_0$ is equivalent to
\begin{equation}\label{eq:uniform}
 B_0\subseteq R[X]\quad\text{for every }X\in\Pos(\cF).
\end{equation}
Necessity is obtained by testing singleton output filters $\uparrow\{y\}$
for $y\in B_0$. For sufficiency, any set positive for $\uparrow B_0$ meets
$B_0$, so it meets $R[X]$. Finally, \eqref{eq:uniform} says that each column
$C_y(R)$, $y\in B_0$, meets every $\cF$-positive set. Lemma~\ref{lem:dual}
identifies this with membership in $\cF$.
\end{proof}

\begin{proposition}\label{prop:oblique}
The family $I_{\ltimes}(\cF,\cG)$ is a filter. If both inputs are proper,
it is proper. In particular
\[
 \Up(\cF\ltimes\cG)=I_{\ltimes}(\cF,\cG).
\]
Moreover, for all filters the non-strict chain corresponding to
\eqref{eq:problem} holds.
\end{proposition}
\begin{proof}
The full relation has witness $B_0=V$, and upward closure follows by enlarging
columns. For two relations with witnesses $B_1,B_2$, use $B_1\cap B_2\in\cG$.
On that witness set, the columns of their intersection belong to $\cF$,
since $C_y(R_1\cap R_2)=C_y(R_1)\cap C_y(R_2)$. This proves intersection
closure. For proper inputs the witness set is nonempty, and every one of its
columns is nonempty; consequently no member of the family is empty.

Every rectangle in $I_{\square}$ belongs to $I_{\ltimes}$, with its right
side as witness. If $R\in I_{\ltimes}$ has witness $B_0$, then every
$\cG$-positive $Y$ meets $B_0$, and every $\cF$-positive $X$ satisfies
$B_0\subseteq R[X]$. Thus $I_{\ltimes}\subseteq I_2$.
Taking generated families and reversing inclusion proves the chain.
\end{proof}

\section{The cofinite filter}

Let $\Phi=\{A\subseteq\N:\N\setminus A\text{ is finite}\}$.
It is proper and non-principal, and $\Pos(\Phi)$ consists exactly of the
infinite subsets of $\N$. Lemmas~\ref{lem:dual} and~\ref{lem:indices} give
\begin{align*}
 R\in I_2(\Phi,\Phi)
 &\Longleftrightarrow R[X]\text{ is cofinite for every infinite }X,\\
 R\in I_{\ltimes}(\Phi,\Phi)
 &\Longleftrightarrow\text{cofinitely many columns of }R\text{ are cofinite}.
\end{align*}

\begin{theorem}\label{thm:nat}
The proper, non-principal input filter $\Phi$ satisfies
\[
 \Phi\times_F^{\RLD}\Phi=\bot
 \sqsubset\Phi\ltimes\Phi
 \sqsubset\Phi\times^{\RLD}\Phi.
\]
\end{theorem}
\begin{proof}
Put $L=\{(x,y)\in\N^2:y<x\}$ and $H=\{(x,y)\in\N^2:x<y\}$.
For infinite $X$, unboundedness gives $L[X]=\N$, whereas
$H[X]=\{y:\min X<y\}$ is cofinite. Thus $L,H\in I_2(\Phi,\Phi)$.
Since $L\cap H=\varnothing$, the second product's generated family contains
the empty relation and is the improper filter. Proposition~\ref{prop:oblique}
makes the oblique product proper, so the first inequality is strict.

Let $D=\{(x,y)\in\N^2:x\ne y\}$. Every column of $D$ is cofinite;
hence $D\in\Up(\Phi\ltimes\Phi)$. On the other hand, every cofinite
rectangle $A\times B$ contains $(n,n)$ for some $n\in A\cap B$.
Thus no such rectangle is contained in $D$. Finite intersections of cofinite
rectangles are again cofinite rectangles, so
$D\notin\Up(\Phi\times^{\RLD}\Phi)$. With the non-strict chain, this
proves strictness of the second inequality.
\end{proof}

\begin{corollary}
Both existential conjectures in Open Problem Garden problem 37339 are true.
\end{corollary}

\section{The proposed real-number example}

On $\R$, let
\[
 \Delta=\left\{S\subseteq\R:
       (-\varepsilon,\varepsilon)\subseteq S\text{ for some }\varepsilon>0
       \right\}.
\]
This is the neighbourhood filter of $0$, equivalently the filter meet of the
principal symmetric intervals in the problem's source. Its meet with the
principal positive half-line is
\begin{equation}\label{eq:delta}
 \Delta_+=\Delta\sqcap\uparrow^{\R}(0,+\infty)
 =\{S\subseteq\R:\exists\varepsilon>0,\ (0,\varepsilon)\subseteq S\}.
\end{equation}
The family in \eqref{eq:delta} is proper because every positive interval is
nonempty. It is non-principal because the intersection of all those intervals
is empty. Its positive sets satisfy
\begin{equation}\label{eq:positive-real}
 X\in\Pos(\Delta_+)
 \quad\Longleftrightarrow\quad
 X\cap(0,\varepsilon)\ne\varnothing\quad\text{for every }\varepsilon>0.
\end{equation}

\begin{theorem}\label{thm:real}
For the proper, non-principal real-number filter in \eqref{eq:delta},
\[
 \Delta_+\times_F^{\RLD}\Delta_+=\bot
 \sqsubset\Delta_+\ltimes\Delta_+
 \sqsubset\Delta_+\times^{\RLD}\Delta_+.
\]
Consequently the particular pair proposed in the problem also has an
affirmative answer.
\end{theorem}
\begin{proof}
Use the disjoint relations $L=\{(x,y)\in\R^2:y<x\}$ and
$H=\{(x,y)\in\R^2:x<y\}$. For a $\Delta_+$-positive set $X$, choose
$x\in X\cap(0,1)$. The image $L[X]$ contains $(0,x)$, a member of
$\Delta_+$. For every $y>0$, choose $x\in X\cap(0,y)$; then $y\in H[X]$.
Thus $H[X]$ contains the positive half-line, also a member of $\Delta_+$.
Both images meet every $\Delta_+$-positive output set, so Lemma~\ref{lem:indices}
gives $L,H\in I_2(\Delta_+,\Delta_+)$. Their disjointness again forces the
second product to be $\bot$, and the oblique product is proper by
Proposition~\ref{prop:oblique}.

For $D=\{(x,y)\in\R^2:x\ne y\}$, every column $\R\setminus\{y\}$
contains a positive interval: use $(0,y)$ when $y>0$ and $(0,1)$ otherwise.
Hence all columns belong to $\Delta_+$, and $D$ belongs to the oblique
family with witness $\R$. Any two members $A,B$ of $\Delta_+$ intersect,
so their rectangle contains a diagonal point and cannot be contained in $D$.
As finite intersections of such rectangles remain rectangles of members,
$D$ is absent from the classical family. This proves the second strict
inequality.
\end{proof}

\begin{remark}
The proof uses only finite intersection closure of filters. The collapse does
not assert that a meet of the corresponding funcoids is the principal funcoid
of the empty relation: it is the meet in the \emph{reloid} lattice that is
computed by intersecting relation generators. Both input filters are proper;
the least reloid is permitted as an output of the second product.
\end{remark}

\section{Formal verification}

The accompanying \texttt{lean/StrictProducts.lean} uses Lean 4.27.0 and a pinned
Mathlib revision \cite{mathlib}. Filters and reloids are represented by
Mathlib's \texttt{Filter} type, with its reverse-inclusion order, so their
families are necessarily nonempty. The development defines the required
funcoids as pairs of filter-valued maps satisfying the meshing equivalence,
constructs the principal and product funcoids, and defines the three reloid
products from their generating families. It proves the reduction lemmas and
the chain from those definitions.

The final statements \texttt{frechet\_solution} and
\texttt{real\_solution} include properness, non-principality, collapse of the
second product, and both strict inequalities. The theorem
\texttt{deltaPlus\_eq\_inf} identifies the real filter with the exact meet in
\eqref{eq:delta}; \texttt{delta\_eq\_iInf} identifies $\Delta$ with the
source's meet of symmetric intervals. The real example uses Mathlib's
\texttt{Real} type, not a rational surrogate. A separate
\texttt{lean/Check.lean} checks the headline statement shapes and their axiom
dependencies. The proof does not require the rest of AGT's composition or
category theory.

\section*{AI disclosure}
The writing of the proof was assisted by DeepSeek 4.1 Flash and OpenAI
GPT-6.1 sol.

\begin{thebibliography}{9}
\raggedright

\bibitem{garden}
V. Porton, \emph{Strict inequalities for products of filters},
Open Problem Garden, problem 37339, posted August 9, 2011.
\url{https://www.openproblemgarden.org/op/strict_inequalities_for_products_of_filters}.
Accessed October 9, 2026.

\bibitem{porton}
V. Porton, \emph{General Topology as Ordered Semigroup Actions}, Volume 1
(also circulated as \emph{Algebraic General Topology}). Source revision
\texttt{89eeff3457d1a91c8946caffbd9ff6d8d39142dc}.
Relevant sections: filters; funcoidal and reloidal products; reloids induced by
a funcoid; and ``Second product. Oblique product''.
\url{https://github.com/vporton/algebraic-general-topology}.

\bibitem{erratum}
V. Porton, \emph{A typo in my math book}, December 1, 2016.
\url{https://math.portonvictor.org/2016/12/01/typo-in-my-math-book/}.
The corresponding correction is source commit
\texttt{bd2992f473603e88d4bd32b314787101557f405c}.

\bibitem{mathlib}
The Mathlib community, \emph{Mathlib}, Lean mathematical library, release
\texttt{v4.27.0}, revision
\texttt{a3a10db0e9d66acbebf76c5e6a135066525ac900}.
\url{https://github.com/leanprover-community/mathlib4}.

\end{thebibliography}

\end{document}
